\documentclass[a4paper,11pt]{ltjsarticle}

\usepackage{lmodern}
\usepackage[T1]{fontenc}
\usepackage{luatexja-fontspec}  % LuaTeX-ja用のフォント設定

% 配置'H'の使用
\usepackage{float}

% 余白
\usepackage[top=20truemm,bottom=20truemm,left=20truemm,right=20truemm]{geometry}
\usepackage{changepage}

%参考文献
\makeatletter
\renewcommand{\refname}{References}
\makeatother

% physics2
\usepackage{physics2}
\usephysicsmodule{ab}   % 括弧サイズの自動調整
\usephysicsmodule{diagmat}  % 対角行列
\usephysicsmodule[showleft=2,showtop=2]{xmat}   % nm行列
% physics2 END

% 量子回路
\usepackage{quantikz}

% 数式環境
\usepackage{amsmath,amssymb}
% ベクトル
\usepackage{bm}
% braket記法
\usepackage{braket}
% 微分演算子
\usepackage{fixdif}
\usepackage{derivative}
% new command
% Landau記号
\newcommand{\order}[1]{\mathcal{O}\ab(#1)}
% 勾配・発散・回転
% 勾配
\DeclareMathOperator{\grad}{\nabla}
% 発散
\DeclareMathOperator{\divergence}{\nabla\cdot}
% \divが「÷」と競合するため再定義, 「÷」は \divisionsymbolで表記する
\let\divisionsymbol\div 
\renewcommand{\div}{\divergence}
% 回転
\DeclareMathOperator{\rot}{\nabla\times}
% 勾配・発散・回転 END
% 実部
\renewcommand{\Re}{\operatorname{Re}}
% 虚部
\renewcommand{\Im}{\operatorname{Im}}
% トレース
\newcommand{\Tr}{\operatorname{Tr}}
% rank
\newcommand{\rank}{\operatorname{rank}}
% new command END

% 化学式
\usepackage[version=4]{mhchem}
% SI単位系
\usepackage{siunitx}
\DeclareSIUnit\angstrom{Å}

% 表
\usepackage{multirow}
% 表の太線
\newcommand{\bhline}{\noalign{\hrule height 1pt}}
\newcommand{\bvline}{\vrule width 1pt}

% 画像
\usepackage[colorlinks=true, allcolors=black]{hyperref}
\usepackage{bookmark}
\usepackage{graphicx}
\usepackage{titlesec}

% セクション見出しを「IV. APPLICATIONS」形式にする
\titleformat{\section}[block]
	{\centering\normalfont\large\bfseries}
	{\Roman{section}.}
	{0.5em}
	{\MakeUppercase}
\titlespacing*{\section}{0pt}{1.2ex plus .4ex minus .2ex}{0.8ex}

% ショートサマリー冒頭情報
\newcommand{\ShortSummaryTitle}{Thesis Template}
\newcommand{\ShortSummaryAuthor}{Taro Keio}
\newcommand{\ShortSummaryAffiliation}{keio university}
\newcommand{\ShortSummaryOverview}{Write a short overview of your work in one or two sentences.}

\begin{document}

\twocolumn[
\begin{center}
{\LARGE\bfseries \ShortSummaryTitle\par}
\vspace{4mm}
{\large \ShortSummaryAuthor\par}
\vspace{1mm}
{\normalsize \ShortSummaryAffiliation\par}
\vspace{1mm}
{\small\textbf{Overview:} \ShortSummaryOverview\par}
\end{center}
\vspace{6mm}
] 

\section{Introduction}
This template is designed for first-time users of academic writing in LaTeX.
In this section, explain the background of your topic and why the problem is important.
Keep the explanation simple and start with definitions of key terms.

\section{Motivation}
Describe what motivated this study and what gap exists in previous work.
You can write one paragraph about practical needs and another about academic interest.
If possible, include a clear research question at the end of this section.
\cite{Lloyd2016}
\section{Result}
Summarize your main findings in a concise and readable way.
State what improved, what was confirmed, or what was newly discovered.
When numbers are available, present representative values and conditions.

\section{Methods}
Explain how the study was conducted so that others can reproduce it.
List assumptions, data sources, experimental setup, and analysis procedure.
If a method has multiple steps, describe them in the order they were performed.


\section{Discussion}
Interpret the results and connect them to your original motivation.
Discuss limitations, possible sources of error, and future extensions.
End with a short conclusion that highlights the contribution of the work.
\bibliographystyle{junsrt}
\bibliography{reference}
\end{document}